\documentclass[10pt,a4paper]{amsart}
\usepackage[margin=19mm]{geometry}
\usepackage{amsmath,amssymb,mathtools}
\usepackage[scr=rsfs,cal=euler]{mathalfa}
\usepackage{xfrac}
\usepackage[unicode,hidelinks,bookmarks=false]{hyperref}
\newcommand{\ie}{i.e.\ }
\newcommand{\cf}{cf.\ }
\newcommand{\resp}{respectively\ }
\newcommand{\Subsection}{\subsection}
\providecommand{\nobreakdash}{\nobreak\hbox{-}}
\setlength{\emergencystretch}{2em}
\multlinegap0pt
\usepackage{amssymb}
\usepackage{mathtools}
\newtheorem{lemma}{Lemma}
\DeclareMathOperator{\Rea}{Re}
\title{Morse index, topology, and ends of minimal surfaces with noncompact free boundary}
\author{M\'arcio Batista, Matheus B. Martins}
\date{Source: arXiv:2609.29723v1 (CC BY 4.0)}
\begin{document}
\maketitle

\noindent\emph{Source and licence.} Morse index, topology, and ends of minimal surfaces with noncompact free boundary. Version arXiv:2609.29723v1 (CC BY 4.0), distributed by arXiv under the Creative Commons Attribution 4.0 International licence, http://creativecommons.org/licenses/by/4.0/, as recorded in the arXiv licence metadata for this exact version and shown on the abstract page https://arxiv.org/abs/2609.29723v1. Full licence notice: \texttt{licenses/arxiv-2609.29723-CC-BY-4.0.txt}.\par\smallskip

\noindent\textit{Editorial scope.} Counted unit: the article's lemma lem:radial-capacity (source0.tex lines 759--773). The article's complete original proof of it is delivered with it (source0.tex lines 774--878, one \texttt{proof} environment of 105 lines). No mathematics is written, repaired or re-derived: the statement and the proof are the article's own lines, in the article's own notation, and no retained line is substituted or re-flowed.  No further source-local context is needed: every cross-reference of every retained line resolves inside the two delivered spans.  Nothing was omitted from any retained span, so the package carries no omission record.  No retained line contains a citation, so the wrapper carries no bibliography and no bibliography entry.  No retained line contains a figure environment, an image-inclusion command or drawing code, so no image is delivered and the package declares no asset dependency.  Editorial formatting only, in addition: an AMS \texttt{amsart} preamble that restates the article's own declarations and macros used by the retained lines, namely the environments \texttt{lemma} on separate counters, together with the standard packages those lines need.  The retained environments print in this document as Lemma 1 in order of appearance, whereas the article prints the same items with the numbering of its own full document.  The complete native source of the article as distributed in the arXiv e-print for this exact version is bundled unchanged as \texttt{sources/211693/source0.tex}.
\begin{lemma}\label{lem:radial-capacity}
On a punctured disk or half-disk with metric $e^{2\lambda}|dz|^2$, let $m\ge1$ and set
\[
a_m(r)=r^{1-2m}\int_{\Theta}e^{-2\lambda(r,\theta)}\,d\theta,
\]
where $\Theta=[0,2\pi]$ or $[0,\pi]$, respectively. If
\begin{equation}\label{eq:capacity-condition}
\int_0^{r_0}\frac{dr}{a_m(r)}=\infty,
\end{equation}
there are radial cutoffs $\zeta_\ell$, zero near the puncture and equal to one outside neighborhoods shrinking to it, such that
\[
\int |\nabla\zeta_\ell|_\gamma^2|\omega|_\gamma^2\,dA_\gamma\longrightarrow0
\]
for every form whose holomorphic representative has pole order at most $m$. In particular, the lower bound $e^\lambda\ge C^{-1}r^{-m}$ implies this conclusion.
\end{lemma}

\begin{proof}
Write $\omega=\Rea(f\,dz)$. Since $f$ has a pole of order at
most $m$, after reducing the coordinate neighborhood if necessary,
\(
|f(z)|\le C_\omega |z|^{-m}.
\)
For a radial function $\zeta=\zeta(r)$, the metric
$\gamma=e^{2\lambda}(dr^2+r^2d\theta^2)$ gives
\[
|\nabla\zeta|_\gamma^2=e^{-2\lambda}|\zeta'(r)|^2,
\qquad
|\omega|_\gamma^2=e^{-2\lambda}|f|^2,
\qquad
dA_\gamma=e^{2\lambda}r\,dr\,d\theta.
\]
So, if $\zeta'$ is supported in
$[\varepsilon,\delta]$, then
\begin{align*}
\int |\nabla\zeta|_\gamma^2|\omega|_\gamma^2\,dA_\gamma
&=\int_\varepsilon^\delta |\zeta'(r)|^2r
  \int_\Theta e^{-2\lambda(r,\theta)}
  |f(re^{i\theta})|^2\,d\theta\,dr\\
&\le C_\omega^2
  \int_\varepsilon^\delta|\zeta'(r)|^2a_m(r)\,dr.
\end{align*}

We next minimize this one-dimensional upper bound. The function
$a_m$ is positive and continuous on every compact subinterval
of $(0,r_0)$. For any absolutely continuous $\zeta$ satisfying
$\zeta(\varepsilon)=0$ and $\zeta(\delta)=1$,
Cauchy--Schwarz yields
\[
1
=\left|\int_\varepsilon^\delta\zeta'(r)\,dr\right|^2
\le
\left(\int_\varepsilon^\delta
|\zeta'(r)|^2a_m(r)\,dr\right)
\left(\int_\varepsilon^\delta a_m(r)^{-1}\,dr\right).
\]
Thus the radial energy is bounded below by
\[
A_{\varepsilon,\delta}^{-1},
\qquad
A_{\varepsilon,\delta}
:=\int_\varepsilon^\delta a_m(r)^{-1}\,dr.
\]
Equality holds for the increasing function
\[
\zeta_{\varepsilon,\delta}(r)
=
\frac{1}{A_{\varepsilon,\delta}}
\int_\varepsilon^r a_m(t)^{-1}\,dt,
\qquad \varepsilon\le r\le\delta,
\]
since its derivative is
$A_{\varepsilon,\delta}^{-1}a_m(r)^{-1}$.

Choose $\delta_\ell\downarrow0$. By
\eqref{eq:capacity-condition}, and because $a_m^{-1}$ is
integrable on compact subintervals of $(0,r_0)$,
\[
\int_0^{\delta_\ell}a_m(r)^{-1}\,dr=\infty.
\]
We may therefore choose $0<\varepsilon_\ell<\delta_\ell$ such that
$A_{\varepsilon_\ell,\delta_\ell}\ge\ell$. Extend
$\zeta_{\varepsilon_\ell,\delta_\ell}$ by zero for
$r\le\varepsilon_\ell$ and by one for $r\ge\delta_\ell$,
and denote the resulting function by $\zeta_\ell$.
It takes values in $[0,1]$ and is Lipschitz: its derivative is
bounded on the compact transition interval and vanishes elsewhere.
Moreover, $\zeta_\ell$ is eventually one on every compact subset
away from the puncture, and
\[
\int |\nabla\zeta_\ell|_\gamma^2|\omega|_\gamma^2\,dA_\gamma
\le \frac{C_\omega^2}{A_{\varepsilon_\ell,\delta_\ell}}
\le \frac{C_\omega^2}{\ell}\longrightarrow0.
\]

Finally, if $e^\lambda\ge C^{-1}r^{-m}$ uniformly in $\theta$,
then
\[
a_m(r)
=r^{1-2m}\int_\Theta e^{-2\lambda(r,\theta)}\,d\theta
\le C^2|\Theta|\,r.
\]
Hence $\int_0^{r_0}a_m(r)^{-1}\,dr=\infty$.
In this case one may also use the explicit logarithmic transition
\[
\zeta_\varepsilon(r)=
\frac{\log(r/\varepsilon^2)}{|\log\varepsilon|},
\qquad
\varepsilon^2\le r\le\varepsilon,
\]
extended by zero and one. Indeed,
\[
\int |\nabla\zeta_\varepsilon|_\gamma^2
|\omega|_\gamma^2\,dA_\gamma
\le
\frac{C_\omega^2C^2|\Theta|}{|\log\varepsilon|^2}
\int_{\varepsilon^2}^{\varepsilon}\frac{dr}{r}
=
\frac{C_\omega^2C^2|\Theta|}{|\log\varepsilon|}
\longrightarrow0.
\]
\end{proof}

\end{document}
